\documentclass[11pt]{article}
\usepackage[margin=1in]{geometry}
\usepackage{hyperref}
\usepackage{booktabs}

\title{Irembo Voice AI Intent Classification Report}
\author{}
\date{18th January 2026}

\begin{document}
\maketitle

\noindent\textbf{Code Repository:} \url{https://github.com/Abdulraqib20/irembo-voice-ai}\\
\textbf{Code and reproducibility:} See scripts for preprocessing, training, and evaluation.

\section*{Approach}
I fine-tuned a multilingual transformer, \href{https://huggingface.co/Davlan/afro-xlmr-base}{Davlan/afro-xlmr-base}, for intent classification. This model was selected because it includes African language coverage and handles morphology and code-switching between Kinyarwanda and English. In production, I use a hybrid routing strategy: the transformer is primary, a rule-based classifier provides a fast baseline and secondary fallback, and a Mixtral-8x7B LLM fallback is used for low-confidence or ambiguous cases. The transformer remains the primary model because it delivers higher macro performance across all intents.

\paragraph{Why this fits Voice AI:}
Voice AI systems face unique NLP constraints that guided model selection:
\begin{itemize}
  \item \textbf{ASR noise tolerance:} Speech-to-text introduces errors (``rendez-vous'' $\rightarrow$ ``rendezvous'', ``pasiporo'' $\rightarrow$ ``passport''). Transformers learn robust representations that generalize across spelling variants, unlike brittle keyword matching.
  \item \textbf{Code-switching:} Rwandan users naturally mix Kinyarwanda and English (``Ndashaka kureba status"). AfroXLMR's multilingual pretraining handles cross-lingual subword sharing.
  \item \textbf{Informal phrasing:} Spoken language differs from written text including contractions, hesitations, incomplete sentences. The model sees diverse training examples including informal patterns.
  \item \textbf{Latency requirements:} Voice interactions expect sub-second responses. The transformer runs in 60--100ms on CPU, acceptable for real-time use.
\end{itemize}

\section*{Data Strategy}

\paragraph{Cleaning noisy voice transcripts}
I prioritized \textbf{minimal normalization}, strictly limiting preprocessing to whitespace collapse and schema validation. I retained ASR artifacts (such as phonetic approximations) to ensure the training data matches the noisy input expected from the production ASR, preventing a "clean train / dirty test" mismatch that often degrades real-world performance.

\paragraph{Handling code-switching}
The dataset contains Kinyarwanda (rw), English (en), and mixed utterances. I preserve code-switched tokens intact---no language-specific filtering. AfroXLMR's subword tokenizer learns shared representations across languages. A lightweight language detector (\texttt{src/services/enhanced\_language\_detector.py}) scores utterances using marker words (e.g., ``ndashaka'', ``igeze'' for rw; ``status'', ``application'' for en) and flags mixed when both scores $>$0.

\paragraph{Handling code-switching}
Given the dataset contains Kinyarwanda (rw), English (en), and mixed utterances, I adopted a strategy that preserves code-switched sequences intact without language-specific filtering. AfroXLMR's subword tokenizer learns shared representations across languages. A lightweight custom language detector scores utterances using marker words (e.g., ``ndashaka'', ``igeze'' for rw; ``status'', ``application'' for en) and flags mixed when both scores $>$0.

\paragraph{Small and imbalanced data}
Faced with a limited and imbalanced dataset of 700 samples (561 train, 70 val, 69 test) across 13 intents. I track intent/language distributions in preprocessing and use \textbf{Macro F1-Score} (not accuracy) to weight rare intents equally. The \textbf{Hybrid Fallback mechanism} (rule-based + LLM) reduces errors on underrepresented cases while new data accumulates.

\paragraph{Data quality improvement over time}
I established an \textbf{Active Learning loop} driven by production telemetry. The system automatically flags predictions with confidence scores below 0.7 for human review and relabeling. This creates a continuous feedback cycle to retrain the model on its hardest edge cases, with all iterations tracked in a versioned registry.


\section*{Evaluation}
I report accuracy, Macro Precision, Macro Recall, and Macro F1. Macro metrics are chosen to weight all intents equally, which is critical for public services. Weighted metrics can hide poor performance on rare intents.

\begin{table}[h]
\centering
\begin{tabular}{lcccc}
\toprule
Model & Accuracy & Macro F1 & Macro Precision & Macro Recall \\
\midrule
Rule-Based Baseline & 88.41\% & 0.8937 & 0.9034 & 0.8857 \\
AfroXLMR base & 95.65\% & 0.9563 & 0.9654 & 0.9526 \\
\bottomrule
\end{tabular}
\caption{Overall test performance.}
\end{table}

\paragraph{Per-language accuracy}
English 90.5\%, code-switched 100.0\%, Kinyarwanda 97.4\%. This demonstrates strong performance on the target language and on mixed utterances.

\paragraph{Robustness and edge-case testing}
I test robustness with the transformer, baseline, and hybrid evaluation scripts and also probe the live API. Empty inputs are rejected by request validation, ambiguous or mixed-language text is routed through confidence-based fallback, and long utterances are safely handled by the configured max sequence length (128).

\paragraph{Common failure modes}
Failure modes are identified from the hybrid evaluation confusion matrix and sample errors, and they are logged at inference with explicit fallback reasons. The main patterns are confusable intents (e.g., requirements vs fees) and out-of-scope queries that drive low confidence; when primary or fallback classifiers fail, \texttt{unknown} is returned and the event is recorded.

\paragraph{Suitability for a citizen-facing system}
Suitability is assessed via macro F1, per-language accuracy, and operational safeguards in the API. Each request is logged with confidence, fallback reason, and latency; monitoring tracks low-confidence rate, fallback rate, P95 latency, and language drift. Health/readiness endpoints support safe rollout, and escalation is available when confidence drops below the defined threshold.

\section*{MLOps and Production Considerations}

\subsection*{Model Versioning and Deployment}

I use a simple model registry to track versions, training configuration, evaluation metrics, and the active model with rollback support. The model is deployed behind a FastAPI service with health/readiness checks and both single and batch inference routes. Docker Containerization is used for reproducible deployment and predictable runtime behavior.

\subsection*{Monitoring in Production}

Inference is logged per request with confidence, latency, language, and fallback reason, and metrics are aggregated over a rolling window. Drift is monitored via language distribution shifts; performance is tracked with low-confidence and fallback rates plus P95 latency. Alerts are triggered when thresholds are exceeded, and logs are retained for offline error analysis.

\subsection*{Real-Time Voice AI Pipeline Integration}

The pipeline is integrated as: ASR produces text, language is detected, the transformer predicts intent with confidence, and low-confidence cases are routed through fallback. Predicted intents are then mapped to downstream service actions. The real-time API supports both single and batch classification for online and offline flows.

\subsection*{Fallback Strategies}

I use a three-tier fallback: transformer first, rule-based second, and an LLM (Mixtral-8x7B) as a last resort for ambiguous cases. Fallback is triggered by confidence thresholds and every fallback decision is logged for review. This balances accuracy, latency, and interpretability in a citizen-facing setting.

\section*{Voice AI Pipeline Architecture}

The implemented flow: \\
\textbf{Utterance Text} $\rightarrow$ \textbf{Language Detection} (en/rw/mixed) $\rightarrow$ \textbf{Intent Classification} (Transformer $\rightarrow$ Rule-based $\rightarrow$ Mixtral-8x7B (LLM) fallback) $\rightarrow$ \textbf{Response}. \\
Each request is logged with confidence, fallback reason, and latency for observability.

\section*{Ethics and Public Sector Considerations}

\paragraph{Bias and fairness risks}
Government services must serve all citizens equiatably. Key risks and mitigations:
\begin{itemize}
  \item \textbf{Language Bias:} Kinyarwanda speakers should not receive worse service than English speakers. I evaluate per-language accuracy separately (Kinyarwanda 97.4\% vs English 90.5\%) and prioritize Kinyarwanda performance.
  \item \textbf{Intent Coverage Bias:} Rare intents (e.g., \texttt{complaint\_or\_support\_ticket}) may be underrepresented. Macro F1 weighs all intents equally, preventing majority-class dominance.
\end{itemize}

\paragraph{Explainability:} Confidence scores and fallback reasons are returned and logged, and low-confidence cases are escalated. Intent definitions include canonical examples to support transparent responses.

\paragraph{Responsible AI:} Uncertain predictions are routed through fallback or escalation, and every decision is logged with a request ID for auditability. The system is designed to augment human service delivery, with monitoring to surface degradation early.

\section*{Tradeoffs and Limitations}

\begin{itemize}
  \item \textbf{Accuracy vs. latency:} Transformer (95.65\% acc, $\sim$60--100ms) vs. rule-based (88.4\% acc, $\sim$16ms). I chose accuracy for citizen-facing correctness, with rule-based as fast fallback.
  \item \textbf{Model size vs. deployment cost:} AfroXLMR-base (278M params) requires $\sim$1GB memory. For edge deployment, distillation to a smaller model would trade accuracy for resource efficiency.
  \item \textbf{External API dependency:} Groq LLM fallback adds cost (\$0.05--0.10 per 1K tokens) and latency ($\sim$200--400ms). For cost-sensitive deployments, disable LLM tier or use local smaller LLM.
  \item \textbf{Dataset size:} 700 samples is small; results should be validated with real production traffic. The feedback loop addresses this over time.
  \item \textbf{ASR error propagation:} The classifier sees ASR output, not original speech. Severe transcription errors (e.g., ``passport'' $\rightarrow$ ``pass port'') may confuse the model. ASR confidence filtering could help.
\end{itemize}

\paragraph{Multilingual expansion strategy}
To extend beyond Kinyarwanda/English:
\begin{enumerate}
  \item \textbf{Add language:} Collect 200+ labeled utterances in target language (e.g., French, Swahili).
  \item \textbf{Update language detector:} Add marker words for new language to \texttt{enhanced\_language\_detector.py}.
  \item \textbf{Fine-tune incrementally:} Continue training from existing checkpoint with mixed old + new data to prevent catastrophic forgetting.
  \item \textbf{Evaluate separately:} Report per-language metrics to ensure new language doesn't degrade existing performance.
\end{enumerate}

\section*{Conclusion}

AfroXLMR-base provides a reliable, multilingual intent classifier for the Irembo Voice AI use case. It performs strongly on Kinyarwanda (97.4\%) and code-switched speech (100\%) and integrates seamlessly into a production pipeline with confidence-based fallbacks (rule-based + Groq Mixtral LLM) and comprehensive monitoring. The system is designed to be fair (macro F1 metrics), transparent (logging all decisions), and resilient (three-tier fallback) for a citizen-facing government services platform.

A working demo is available via Docker containerization with FastAPI inference endpoints, demonstrating production readiness with health checks, batch processing, and real-time classification.

\end{document}
